\documentclass[10pt]{article}
\usepackage[margin=0.55in]{geometry}
\usepackage{amsmath,amssymb,booktabs,graphicx,microtype}
\usepackage[hidelinks]{hyperref}
\setlength{\parindent}{0pt}
\setlength{\parskip}{3pt}
\setlength{\abovedisplayskip}{4pt}
\setlength{\belowdisplayskip}{4pt}
\setlength{\textfloatsep}{5pt}
\renewcommand{\arraystretch}{1.08}

\title{\vspace{-1.15cm}\textbf{Two-Block Phase Diagram:\\Supplementary Experiments and Seed Dependence}}
\author{}
\date{\vspace{-1.25cm}}

\begin{document}
\maketitle
\vspace{-0.72cm}

\textbf{Purpose.}
This supplement records the training-history and numerical-control experiments omitted from
the concise theory comparison.  The full sweep contains three independent initialization
seeds (0, 1, 2), 72 decay cells per seed, and two matched protocols per cell.

\textbf{The two protocols were both run.}
\emph{Joint} training starts from the random initialization and trains both degree-2 parents
and the degree-4 target simultaneously, with every parameter trainable.  \emph{Parent-warm}
starts from the same seeded base model, first trains only $E,V_1,W_{d2}$ on the two parent
targets read from $h_1$, and then unfreezes the entire model and trains the identical final
objective as joint.  There are 216 runs of each protocol (432 total).  Thus protocol
differences isolate training history, not architecture or final loss.

\textbf{History changes the circuit.}
The two protocols selected different phase labels in 128 of 216 matched
$(\text{seed},\lambda_0,\lambda_1/\lambda_2)$ comparisons (59.3\%).  Of those 128 changes,
122 were $H_D\!\to H_C$ from joint to parent-warm, three were
suppressed$\to H_C$, and three were suppressed$\to H_D$.  Parent warming therefore exposes
a broad compositional basin that joint training almost always misses.  It is not universally
the lower-objective basin: parent-warm had lower final objective in 90/216 comparisons
(41.7\%), although its mean objective difference was $-0.0488$.

\begin{figure}[h]
\centering
\includegraphics[width=0.94\textwidth]{../analysis_output/two_block_phase_diagram/supplementary/protocol_phase_comparison.pdf}
\vspace{-7pt}
\caption{\textbf{Protocol-level phase diagrams.}  Top: majority phase over three seeds for
joint, parent-warm, and the lower-objective protocol used by the main report.  Bottom:
matched-seed disagreement, parent-warm minus joint objective, and the increase in total
$2+2\!\to4$ share.  The near-unity disagreement band is the compositional accessibility
window opened by parent pretraining.}
\end{figure}

\clearpage
\textbf{Degree-4-only objective comparison.}
To separate target fitting from the parent losses and shared hidden-weight costs, define
\[
L_4=\frac{\gamma_4}{2}\operatorname{MSE}_4,
\qquad
J_4^{\rm spec}=L_4+\frac{\lambda_4}{2}\|w_4\|^2.
\]
The first quantity is the degree-4 data term alone; the second adds the only regularizer
uniquely attributable to the degree-4 target.  Hidden-weight decay cannot be assigned
unambiguously to one task and is therefore excluded.  Here $\gamma_4=2$, so $L_4$ equals
the reported degree-4 MSE.

\begin{figure}[h]
\centering
\includegraphics[width=0.82\textwidth]{../analysis_output/two_block_phase_diagram/supplementary/d4_only_objective_difference.pdf}
\vspace{-7pt}
\caption{\textbf{Degree-4-only parent-warm minus joint differences, averaged over the
three matched seeds.}  Blue means parent-warm has the smaller value and red means joint
does.  The signed-log color normalization is linear inside $|\Delta|<10^{-7}$ and
logarithmic outside it, allowing differences from $10^{-9}$ to order one to remain visible.
Across most active cells both protocols fit parity nearly perfectly; parent-warm usually has
a slightly larger $L_4$ even when it forms the desired circuit.  The isolated order-one blue
cells arise when joint training suppresses the degree-4 task while parent warming keeps it
active.}
\end{figure}

This comparison also explains why the lower-total-objective map cannot be inferred from
degree-4 fit alone: its selection additionally includes the two parent losses and all five
weight-decay terms.  Circuit identity and the last few digits of degree-4 MSE are distinct
diagnostics.  Parent-warm has smaller $L_4$ in only 8/216 matched runs and smaller
$J_4^{\rm spec}$ in 70/216.  Nevertheless, the mean differences are negative because the
few order-one joint-suppression failures outweigh the many positive differences of only
$10^{-9}$--$10^{-3}$.

\clearpage
\textbf{Seed-by-seed robustness.}

\begin{center}
\begin{tabular}{@{}clll@{}}
\toprule
Seed & Joint: $H_C/H_D$/suppressed & Parent-warm: $H_C/H_D$/suppressed & Disagreement \\
\midrule
0 & $3/52/17$ & $46/11/15$ & $44/72$ (61.1\%) \\
1 & $3/52/17$ & $44/13/15$ & $42/72$ (58.3\%) \\
2 & $3/52/17$ & $44/13/15$ & $42/72$ (58.3\%) \\
\midrule
All & $9/156/51$ & $134/37/45$ & $128/216$ (59.3\%) \\
\bottomrule
\end{tabular}
\end{center}

The joint maps have identical phase counts for all three seeds.  Parent-warm seeds 1 and 2
are identical in phase count, while seed 0 differs by only two cells.  The shape of the
compositional window is therefore reproducible; the large protocol effect is not caused by a
single favorable seed.  Seed variation is concentrated near the staircase boundary between
$H_C$ and $H_D$ and near task suppression.

\begin{figure}[h]
\centering
\includegraphics[width=0.70\textwidth]{../analysis_output/two_block_phase_diagram/supplementary/seed_protocol_phase_grids.pdf}
\vspace{-7pt}
\caption{\textbf{All six empirical phase maps.}  Each row fixes the random initialization;
the left and right panels differ only by parent pretraining before the common fully unfrozen
objective.  The repeated left/right contrast makes the history dependence directly visible.}
\end{figure}

\clearpage
\textbf{Numerical control 1: failed uniform-learning-rate launch.}
The first width-1024 launch used float64 full-batch SGD with momentum 0.9 and learning rate
$0.01$ for every parameter group.  This is incompatible with the muP readout initialization
$W=O(1/N)$: individual readout gradients are $O(1)$, so the first updates overwhelm their
initial scale.  Parent-pretraining states became non-finite and every worker raised a
non-finite main objective at step 8.  No valid phase cell was produced.  This establishes a
numerical requirement---scale the readout learning rate as $1/N$---not a circuit-phase result.

\textbf{Numerical control 2: corrected width-1024 smoke test.}
Before restarting the sweep we tested seed 0 at $\lambda_0=0.01$ and
$\lambda_1/\lambda_2=1$, using $\eta_{E,V}=0.01$, $\eta_W=0.01/1024$, 1,000 parent-warmup
steps, and 1,000 final-objective steps.

\begin{center}
\begin{tabular}{@{}lrrrrrr@{}}
\toprule
Protocol & Phase & Objective & $A_4$ & $R$ & $f_C$ & $\|\nabla J\|$ \\
\midrule
Joint & $H_D$ & 2.29699 & 0.9453 & 0.4454 & 0.1131 & 1.0490 \\
Parent-warm & $H_C$ & 2.28398 & 0.9445 & 0.5879 & 0.5844 & 0.9698 \\
\bottomrule
\end{tabular}
\end{center}

Both runs stayed finite, and the parent-warm run already entered the compositional basin.
Because the gradient norms were still order one, this was only a stability and basin test;
it was not used as stationary phase evidence.

\textbf{Finite-time qualification.}
Only 99/432 full-sweep runs met the plateau criterion; 333 reached the 50k-step cap.  The
phase maps therefore describe circuits found by the stated finite-time protocols, while the
lower-objective aggregation is only a best-observed proxy for the zero-temperature theory.

\begin{figure}[h]
\centering
\includegraphics[width=0.68\textwidth]{../analysis_output/two_block_phase_diagram/supplementary/convergence_diagnostics.pdf}
\vspace{-7pt}
\caption{\textbf{Convergence diagnostics.}  Mean training fraction and final raw gradient
norm over the three seeds.  Weak and moderate decay usually consume the full budget; rapid
termination is concentrated in strongly regularized or suppressed cells.}
\end{figure}

\end{document}
